\documentclass[12pt,a4paper]{article}

\usepackage[utf8]{inputenc}
\usepackage[T1]{fontenc}
\usepackage[portuguese]{babel}

\usepackage[colorlinks=true, urlcolor=blue, linkcolor=black]{hyperref}
\usepackage{geometry}

\geometry{
    a4paper,
    left=3cm,
    right=3cm,
    top=3cm,
    bottom=3cm,
}

\begin{document}

\begin{center}
    {\Large \textbf{Índice - LoRa (Long Range) \& LoRaWAN}} \\[0.6cm]
    {\large \textbf{Danilo Davi Gomes Fróes}} \\[0.4cm]
    {\large \textbf{Grupo 18}}
\end{center}

\vspace{1cm}

\begin{enumerate}
    \item Introdução
        \begin{itemize}
            \item IoT
            \item LPWAN
            \item Diferença: LoRa e LoRaWAN
        \end{itemize}

    \item Modulação LoRa (Camada Física)
        \begin{itemize}
            \item Chirp Spread Spectrum (CSS)
            \item Parâmetros de Rádio-Frequência
            \item Link Budget
        \end{itemize}

    \item Protocolo LoRaWAN (Camada de Enlace e Rede)
        \begin{itemize}
            \item Arquitetura
            \item Classes de Dispositivos
            \item Network Server
        \end{itemize}

    \item Segurança em LoRaWAN
        \begin{itemize}
            \item Criptografia Ponta-a-Ponta
            \item NwkSKey e AppSKey (Chaves de Sessão)
            \item OTAA e ABP
        \end{itemize}

    \item Análise Comparativa
        \begin{itemize}
            \item LoRaWAN vs NB-IoT e Sigfox
            \item LoRaWAN vs Redes Locais
        \end{itemize}

    \item Implementação (Exemplo de Uso Prático)
        \begin{itemize}
            \item Circuito e Componentes
            \item Código
            \item Testes da Implementação
        \end{itemize}

    \item Conclusão

    \item Referências Bibliográficas
\end{enumerate}

\end{document}